\documentclass[12pt,landscape,a4paper]{article}
\usepackage[utf8]{inputenc}
\usepackage[T1]{fontenc}
\usepackage[spanish]{babel}
\usepackage{amsmath,amssymb}
\usepackage{graphicx}
\usepackage{geometry}
\usepackage{enumitem}
\usepackage[colorlinks=true, citecolor=blue, linkcolor=blue]{hyperref}
\usepackage{xcolor}

\geometry{margin=0.7cm}
\pagestyle{empty}
\sloppy

\definecolor{azul-ucv}{RGB}{0, 51, 102}
\definecolor{rojo-accento}{RGB}{180, 30, 30}

\newcommand{\fila}[3]{%
  \noindent
  \begin{minipage}[t]{\textwidth}
    \noindent
    \begin{minipage}[t]{0.38\textwidth}
      \centering
      \vspace{0pt}
      \includegraphics[width=0.95\textwidth, height=3.2cm, keepaspectratio]{#3}
    \end{minipage}
    \hfill
    \begin{minipage}[t]{0.58\textwidth}
      \vspace{0pt}
      \textbf{\textcolor{azul-ucv}{\large #1}}\\[2pt]
        \small #2
    \end{minipage}
  \end{minipage}
  \vspace{0.1cm}
}

\begin{document}

% Encabezado
\begin{center}
  \textcolor{azul-ucv}{\Large\textbf{Simbología de Instrumentos Analógicos}}\\[1pt]
  \textcolor{rojo-accento}{\small Guía de referencia rápida para mediciones eléctricas}\\[2pt]
  \rule{\textwidth}{1pt}
\end{center}

\vspace{0.1cm}

% ===== A =====
\fila{A. Magnitud Eléctrica}
  {Es la única que se imprime en forma independiente y sobre cualquier ángulo de la esfera. Indica la magnitud que mide el instrumento: intensidad (A), tensión (V), resistencia ($\Omega$), potencia (W), frecuencia (Hz), etc.}
  {img/magnitud-electrica.jpeg}

\noindent\rule{\textwidth}{0.4pt}

% ===== B =====
\fila{B. Clase de Corriente}
  {Indica el tipo de corriente para el cual está diseñado el instrumento:
  \begin{itemize}[nosep, leftmargin=*, before=\vspace{1pt}, after=\vspace{1pt}]
    \item[] \textbf{——} Corriente continua (CC)
    \item[] $\sim$ Corriente alterna (CA)
    \item[] $\cong$ Corriente continua y alterna (CC/CA)
    \item[] $\approx$ Corriente alterna trifásica
  \end{itemize}}
  {img/clase-corriente.jpeg}

\noindent\rule{\textwidth}{0.4pt}

% ===== C =====
\fila{C. Seguridad en la Manipulación}
  {Las fallas de aislamiento provocan diferencia de potencial entre la caja y tierra. Los aparatos se someten a una tensión de prueba (\textit{tensión de ensayo}) indicada con una estrella con un número en su interior, representando la tensión en kilovoltios (kV).}
  {img/tabla-tension-ensayo.jpeg}

% ===== D =====
\fila{D. Posición de Funcionamiento}
  {Indica la orientación en la que debe operar el instrumento para garantizar su precisión:
  \begin{itemize}[nosep, leftmargin=*, before=\vspace{1pt}, after=\vspace{1pt}]
    \item[] \rule{1em}{0.5pt} Horizontal
    \item[] $\uparrow$ Vertical
    \item[] $\angle\,60^\circ$ Inclinada
  \end{itemize}}
  {img/posicion-funcionamiento.jpeg}

\noindent\rule{\textwidth}{0.4pt}

% ===== E =====
\fila{E. Clase de Precisión}
  {Número que indica el error porcentual del instrumento. La precisión se expresa en \% del valor a plena escala.
  \begin{itemize}[nosep, leftmargin=*, before=\vspace{1pt}, after=\vspace{1pt}]
    \item[] \textbf{Precisión:} 0.1, 0.2, 0.5
    \item[] \textbf{Industriales:} 1, 1.5, 2.5, 5
  \end{itemize}
  $\displaystyle \varepsilon = \frac{\text{clase}}{100} \times A_{\text{máx}}$}
  {img/clase-precision.jpeg}

\noindent\rule{\textwidth}{0.4pt}

% ===== F =====
\fila{F. Mecanismo de Funcionamiento}
  {Mediante símbolos gráficos se distingue el principio de operación interno del instrumento:
  \begin{itemize}[nosep, leftmargin=*, before=\vspace{1pt}, after=\vspace{1pt}]
    \item[] Bobina móvil (imán permanente)
    \item[] Hierro móvil
    \item[] Electrodinámico (bobinas fija y móvil)
    \item[] Rectificador
  \end{itemize}}
  {img/mecanismo-funcionamiento.jpeg}

\vfill

% Pie de página
\noindent\rule{\textwidth}{0.4pt}\\[2pt]
\noindent
{\small
  \textcolor{gray}{Basado en:} J.~Blondell, \textit{Simbología de los Instrumentos Analógicos para Mediciones Eléctricas}. \\
  \textcolor{gray}{Prof.} Emerson Warhman \hfill \textcolor{gray}{\today}
}

\end{document}
